\begin{definition}[{\cite[Definition 2.4.]{clausen_scholze_cmcg}}] \label{definition:kappa_condensed_set_for_a_cardinal}
    Let $\kappa$ be a \CrefAndHyperrefIfExist{definition:cardinal_number_and_cardinality_of_a_set}{cardinal}. 
    A \hldef{$\kappa$-condensed set} is a \CrefAndHyperrefIfExist{definition:sheaf_on_a_site}{sheaf of sets} on \CrefAndHyperrefIfExist{definition:kappa_small_profinite_set_for_a_cardinal}{$\operatorname{Prof}_{\kappa}$}.

    The category of $\kappa$-condensed sets is denoted by \hl{$\operatorname{CondSet}_{\kappa}$}.

    Explicitly, a $\kappa$-condensed set $X$ is a \CrefAndHyperrefIfExist{definition:functor_between_categories}{functor}
    $$ X: \operatorname{Prof}_{\kappa}^{o p} \rightarrow \text { Set } $$
    such that

    \begin{enumerate}
        \item If $\left(S_{i}\right)_{i \in I}$ is a finite collection of objects of $\operatorname{Prof}_{\kappa}^{o p}$, then

        $$ X\left(\sqcup_{i \in I} S_{i}\right) \xrightarrow{\sim} \prod_{i \in I} X\left(S_{i}\right) ; $$

        \item If $f: T \rightarrow S$ is a surjective map in $\operatorname{Prof}_{\kappa}$, then

        $$ X(S) \stackrel{f^{*}}{\hookrightarrow} X(T) $$

        is injective with image equal to the set of those $x \in X(T)$ for which $p_{1}^{*} x=p_{2}^{*} x$, where $p_{1}, p_{2}: T \times_{S} T \rightarrow T$ are the two projections.
    \end{enumerate}


% There are two ways of thinking about this definition. The first is that instead of encoding a topological space structure on $X$ directly, we encode, for every $S \in$ Prof, the data of an abstract set $X(S)$ which we think of as specifying the set of "continuous maps from $S$ to $X$ ". From this perspective, the above structure and conditions should make intuitive sense:

% (1) the contravariant functoriality $X: \operatorname{Prof}_{\kappa}^{o p} \rightarrow$ Set comes from the idea that a "continuous map" $S \rightarrow X$ and a continuous map $T \rightarrow S$ should compose to give a "continuous map" $T \rightarrow X$

% (2) the first axiom encodes that giving a "continuous map" out of a disjoint union should be the same as separately giving a "continuous map" out of each constituent piece;

% (3) the second axiom encodes that if $f: T \rightarrow S$, then giving a "continuous map" out of $S$ should be the same as giving a "continuous map" out of $T$ which is constant on the fibers of $f$.

% To make this last point more palatable, note that a continuous surjection between compact Hausdorff spaces is a topological quotient map.

% The other way of think about this definition, actually equivalent to the first though perhaps psychologically different, comes back to the idea that profinite sets are the building blocks for arbitrary condensed sets. More formally, the Yoneda embedding gives a fully faithful functor

% $$ \operatorname{Prof}_{\kappa} \hookrightarrow \text { CondSet }_{\kappa} $$ 

% defined by

% $$ S \mapsto(\underline{S}: T \mapsto \operatorname{Cont}(T, S)) $$

% allowing us to view $\kappa$-small profinite sets as special kinds of $\kappa$-condensed sets; and an arbitrary $\kappa$-condensed set is built from $\kappa$-small profinite sets via colimits in CondSet ${ }_{\kappa}$.

% From this perspective, the role of the sheaf conditions is perhaps a bit more tangible. If we had used the presheaf category instead of sheaf category in the definition, then CondSet $_{\kappa}$ would be freely generated under colimits by $\operatorname{Prof}_{\kappa}$, so every relation between colimits would be essentially diagrammatic in nature. The role of the sheaf conditions is to enforce that certain colimits we already have in $\operatorname{Prof}_{\kappa}$ should be preserved by the Yoneda embedding, and thus be visible in CondSet C $_{\kappa}$. Indeed:

% (1) the first sheaf condition says that given finitely many profinite sets $S_{i}$, we have

% $$ \sqcup_{i \in I} \underline{S_{i}} \xrightarrow{\sim} \underline{\sqcup_{i} S_{i}} $$

% (2) the second sheaf condition says that that if $T \rightarrow S$, thereby presenting $S$ as the quotient of $T$ by the equivalence relation $T \times_{S} T$, then $\underline{S}$ is the quotient of $\underline{T}$ by the equivalence relation $\underline{T \times_{S} T}$.

% Carrying around the cardinal $\kappa$ is sometimes a bit annoying. If $\kappa< \kappa^{\prime}$, then there is a pullback functor

% $$ \pi^{*}: \text { CondSet }_{\kappa} \rightarrow \text { CondSet }_{\kappa^{\prime}} $$

% characterized by the fact that it preserves colimits and restricts to the obvious inclusion Prof ${ }_{\kappa} \rightarrow$ Prof $_{\kappa^{\prime}}$ under the Yoneda embeddings. Now it turns out that, for $\kappa$ and $\kappa^{\prime}$ restricted to a cofinal class of cardinals (e.g., strong limit cardinals of sufficiently large cofinality), the functor $\pi^{*}$ is fully faithful and preserves limits and internal hom's of any fixed size. This means that the colimit category

% $$ \text { CondSet }:=\underset{\kappa}{\varinjlim } \text { CondSet }_{\kappa} $$

% of condensed sets is just as well-behaved as each individual CondSet ${ }_{\kappa}$ in terms of its topos-theoretic exactness properties, because calculations can always be done at some "finite" level. Moreover, CondSet is generated under colimits by arbitrary profinite sets in the same way CondSet $_{\kappa}$ is generated under colimits by $\kappa$-small profinite sets.

% Let's try to understand a bit about what condensed sets can look like by carving out a hierarchy of full subcategories of CondSet:

% $$ \text { qcProj } \subset \text { Prof } \subset \text { qcqs } \subset \text { qs } \subset \text { CondSet } $$

% We already know about Prof: these were our basic building blocks, the profinite sets. The other three full subcategories are given by some general properties an object in a topos can have. We review the relevant definitions in the current context.
\end{definition}